\begin{textsample}{POS Dim 3 – rn – Score 28.00 – rn/rn\_inf\_32.txt}  \label{ex:f3_pos_092}
APROFUNDANDO O TEMA \textbf{Sugestões} de \textbf{Brinquedos} e Materiais para Educação \textbf{Infantil} - \textbf{Bebês} ( 0 a 1 ano e meio ) Chocalhos, móbiles sonoros, sinos, \textbf{brinquedos} para morder, bolas de 40 cm e \textbf{menores}, blocos macios, \textbf{livros} e imagens coloridos, \textbf{brinquedos} de empilhar, encaixar, espelhos.
Objetos com diferentes \textbf{texturas} ( mole, rugoso, liso, duro ) e coloridos, que fazem \textbf{som} ( \textbf{brinquedos} musicais ou que \textbf{emitem} \textbf{som} ), de movimento ( \textbf{carros} e objetos para empurrar ), para encher e esvaziar. \textbf{Brinquedos} de \textbf{parque}. \textbf{Brinquedos} para \textbf{bater}.
Cesto com objetos de materiais naturais, metal e de uso cotidiano.
Colcha, rede e colchonete.
Bichinhos de pelúcia.
Estruturas com blocos de espuma para subir, descer, entrar em túneis. - \textbf{Crianças} \textbf{pequenas} ( 1 ano e meio a 3 anos e 11 \textbf{meses} ) Túneis, caixas e espaços para entrar e \textbf{esconder} -se, \textbf{brinquedos} para empurrar, puxar, bolas, quebra-cabeças simples, \textbf{brinquedos} de \textbf{bater}, \textbf{livros} de história, fantoches e teatro, blocos, encaixes, jogos de memória e de percurso, animais de pelúcia, bonecos/as, massinha e tinturas de \textbf{dedo}.
Bonecas/os, \textbf{brinquedos}, \textbf{mobiliário} e acessórios para o faz de conta.
Sucata \textbf{doméstica} e industrial e materiais da natureza.
Sacolas e latas com objetos diversos de uso cotidiano para exploração.
TV, computador, aparelho de \textbf{som}, CD.
Triciclos e \textbf{carrinhos} para empurrar e dirigir.
Tanques de \textbf{areia}, \textbf{brinquedos} de \textbf{areia} e água, estruturas para trepar, subir, descer, balançar, \textbf{esconder}.
Bola, corda, bambolê, papagaio, perna de pau, amarelinha.
Materiais de artes e construções. \textbf{Tecidos} diversos.
Bandinha rítmica. - \textbf{Crianças} Maiores Pré-escolares ( 4 e 5 anos e 11 \textbf{meses} ) Boliches, jogos de percurso, memória, quebra-cabeça, dominó, blocos lógicos, loto, jogos de profissões e com outros temas.
Materiais de arte, \textbf{pintura}, desenho.
CD com músicas, danças.
Jogos de construção, \textbf{brinquedos} para faz de conta e acessórios para \textbf{brincar}, teatro e fantoches.
Materiais e \textbf{brinquedos} estruturados e não estruturados.
Bandinha rítmica. \textbf{Brinquedos} de \textbf{parque}.
Tanques de \textbf{areia} e materiais diversos para \textbf{brincadeiras} na água e \textbf{areia}.
Sucata \textbf{doméstica} e industrial, materiais da natureza.
Papéis, papelão, cartonados, revistas, jornais, gibis, cartazes e folhas de propaganda.
Bola, corda, bambolê, pião, papagaio, 5 marias, bilboquê, perna de pau, amarelinha, varetas gigantes.
Triciclos, \textbf{carrinhos}, equipamentos de \textbf{parque}. \textbf{Livros} \textbf{infantis}, letras móveis, material dourado, globo, mapas, lupas, balança, peneiras, copinhos e \textbf{colheres} de medida, gravador, TV, máquina \textbf{fotográfica}, aparelho de \textbf{som}, computador, impressora. ( KISHIMOTO, 2010, pp. 17-18 ) As produções expostas nos espaços, internos ou externos, são cheias de significados que direcionam e educam o olhar, e ainda podem ser referência para o pensamento das \textbf{crianças}.
As \textbf{crianças} podem se ver nos \textbf{registros} expostos e acompanhar seus processos de aprendizagem.
Além das salas de referências e dos espaços de circulação e de encontros coletivos cobertos ou ao ar livre, na medida do possível, pode -se organizar ambientes específicos para determinadas atividades educativas, como salas de leitura e de multimídia, espaços para apresentações culturais de teatro, dança ou saraus literários e musicais, atelier de artes e \textbf{parques}.
Cada \textbf{instituição} organiza seus ambientes de aprendizagem conforme as possibilidades de espaços físicos que dispõem, no entanto, podem ser pensadas alternativas de garantir materiais diversificados as \textbf{crianças} mesmo em \textbf{pequenos} espaços, como nos exemplos de \textbf{bibliotecas} itinerantes, baú da fantasia, caixas com jogos de construção e materiais alternativos, espaços coletivos para \textbf{brinquedos}, entre outros.
Os corredores podem se transformar em ambientes de interações das \textbf{crianças} com a disposição de artefatos para experimentar sentidos e expressões.
Este documento orienta as possibilidades de organização do tempo, espaços, materiais e relações aqui elencadas e abre caminho para que cada \textbf{instituição}, considerando suas \textbf{crianças} e sua proposta pedagógica, garanta boas condições de interações e aprendizagens.

% matched lemmas: areia, bater, bebê, biblioteca, brincadeira, brincar, brinquedo, carro, colher, criança, dedo, doméstico, emitir, esconder, fotográfico, infantil, instituição, livro, mobiliário, mês, parque, pequeno, pintura, registro, som, sugestão, tecido, textura
\end{textsample}
